% Contenido del Anexo 2 - Marco Teórico

\subsection{Introducción}

El presente anexo desarrolla una versión ampliada del Marco Teórico asociado al proyecto de desarrollo del sistema de gestión de análisis y muestras para el Instituto Nacional de Investigación Agropecuaria (INIA). El objetivo es ofrecer una mirada más completa sobre el contexto en el que se inscribe el proyecto: qué soluciones existían antes, por qué se necesitaba un cambio y cuáles son las tendencias actuales en digitalización dentro del ámbito científico y administrativo.

A medida que las organizaciones comienzan a depender más de la información que generan, las herramientas que antes parecían suficientes, como las planillas Excel o los registros manuales, empezaron a quedarse cortas frente a la complejidad y el volumen de información que debe manejarse hoy. Esto se da, especialmente, en los entornos de investigación, donde contar con datos precisos, trazables y siempre disponibles es esencial.

Esta es la principal razón por la cual cada vez más instituciones están migrando hacia sistemas web que centralizan la información, la organizan de forma más segura y permiten trabajar con reglas claras, validaciones y controles. Este documento profundiza en ese proceso, retomando conceptos teóricos, referencias tecnológicas y comparaciones con herramientas existentes para explicar por qué el desarrollo de un sistema a medida resulta adecuado para el caso del INIA.

El propósito de este capitulo es servir como base para entender por qué INIA necesita un sistema más moderno y por qué las soluciones tradicionales ya no alcanzan. Aquí se describen conceptos y antecedentes que explican la transición de los métodos manuales hacia plataformas digitales más robustas.

\subsubsection{Limitaciones de las hojas de cálculo}

En estos últimos años las organizaciones adoptaron hojas de cálculo como herramienta base para almacenar información gracias a su accesibilidad, bajo costo, flexibilidad y facilidad de uso. Sin embargo, a medida que los procesos se vuelven más complejos y los volúmenes de datos crecen, estos mecanismos se vuelven insuficientes y arriesgados para labores críticas. Las limitaciones más comunes incluyen:

\begin{itemize}
    \item \textbf{Escalabilidad restringida:} Las hojas de cálculo no están diseñadas para manejar grandes volúmenes de datos ni crecer de forma sostenible. A medida que aumentan los registros, las operaciones se vuelven lentas, propensas a fallos y difíciles de mantener.
    
    \item \textbf{Colaboración limitada:} La funcionalidad de edición simultánea que es ofrecida por la mayoría de las aplicaciones de gestión de hojas de cálculo es vulnerable a conflictos, sobreescrituras y pérdida de información.
    
    \item \textbf{Falta de trazabilidad y mecanismos de auditoría:} Resulta difícil rastrear cambios, identificar responsables y asegurar integridad de datos.
    
    \item \textbf{Integración deficiente con sistemas externos:} La conexión con APIs, bases de datos, servicios externos u otros sistemas institucionales resulta limitada.
    
    \item \textbf{Fragmentación de la información:} Es habitual que diferentes usuarios creen y mantengan múltiples versiones de una misma planilla, produciendo duplicados, inconsistencias y pérdida de coherencia en los datos.
    
    \item \textbf{Dependencia del conocimiento del creador:} Muchas planillas contienen fórmulas o lógicas internas que solo entiende quien las creó; si esa persona no está disponible, su mantenimiento se vuelve complejo o directamente inviable.
    
    \item \textbf{Riesgo de pérdida o daño:} Cuando los archivos se guardan localmente, la información queda expuesta a fallos de hardware, extravío o eliminación accidental, sin mecanismos de recuperación robustos.
    
    \item \textbf{Falta de estandarización:} Cada usuario puede modificar formatos, columnas o fórmulas sin ningún tipo de control, generando desorden y dificultando el trabajo conjunto.
    
    \item \textbf{Ausencia de flujos de trabajo formales:} Las planillas no permiten establecer procesos, roles, estados o validaciones complejas, lo que limita la organización y vuelve más difícil garantizar que las tareas se realicen de forma consistente.
\end{itemize}

\subsubsection{La Transición Hacia Sistemas Web}

Las planillas, aunque al inicio eran muy útiles para las organizaciones, no alcanzaban para sostener procesos que crecían en volumen, complejidad y necesidad de control. A medida que la información aumentaba y los equipos debían trabajar de forma más coordinada, se volvió habitual buscar alternativas más sólidas que permitieran centralizar datos, automatizar tareas y reducir la dependencia del manejo manual.

En este contexto surgieron los sistemas web como su evolución. Estas plataformas permiten que toda la información se concentre en un único entorno, accesible desde distintos dispositivos y ubicaciones, y con reglas claras que garantizan organización y coherencia. Además, reemplazan tareas manuales propensas a errores por mecanismos automáticos que validan, controlan y registran cada acción realizada.

Los sistemas modernos ofrecen una serie de ventajas que los convierten en la opción más eficiente cuando se requiere confiabilidad y estructura:

\begin{itemize}
    \item \textbf{Acceso multiplataforma:} es posible trabajar desde computadoras, tablets o teléfonos, sin necesidad de instalar software adicional.
    
    \item \textbf{Escalabilidad horizontal:} gracias a arquitecturas distribuidas, los sistemas pueden crecer sin perder rendimiento, incluso frente a un aumento importante de datos o usuarios.
    
    \item \textbf{Integración nativa mediante APIs REST:} los sistemas pueden comunicarse entre sí, permitiendo automatizar procesos, importar y exportar información o conectar servicios externos.
    
    \item \textbf{Validaciones automáticas a nivel de negocio:} reglas que antes dependían de cada usuario ahora se aplican de forma consistente desde el sistema, reduciendo errores y asegurando calidad de los datos.
    
    \item \textbf{Auditorías completas y seguimiento de operaciones:} cada acción queda registrada, facilitando el control, la trazabilidad y la revisión de procesos.
    
    \item \textbf{Estandarización de flujos de trabajo:} se definen pasos claros, roles, permisos y estados, garantizando que todos los usuarios sigan el mismo procedimiento.
\end{itemize}

Por estas razones, este tipo de soluciones se ha vuelto el estándar en laboratorios, centros de investigación y organizaciones científicas. En estos entornos es fundamental asegurar la trazabilidad de muestras, el control sobre los análisis y la consistencia de los datos. Los sistemas web permiten cumplir con estas necesidades de forma confiable y escalable, convirtiéndose en la base tecnológica sobre la cual muchas instituciones organizan hoy sus procesos técnicos y administrativos.

\subsubsection{Sistemas LIMS como referencia conceptual}

Los LIMS (Laboratory Information Management Systems) son sistemas especialmente diseñados para gestionar muestras y procesos de laboratorio. Aunque el sistema del INIA no es un LIMS completo, adopta varias de sus ideas, como:

\begin{itemize}
    \item Organización clara de fichas y registros.
    \item Trazabilidad completa del ciclo de vida de una muestra.
    \item Gestión por roles con permisos diferenciados.
    \item Validaciones específicas basadas en reglas técnicas.
    \item Escalabilidad para manejar altos volúmenes de datos.
\end{itemize}

Estos conceptos sirvieron como guía para definir la estructura del proyecto.

\subsection{Tecnologías Relevantes}

El proyecto se desarrolló utilizando un conjunto de tecnologías modernas, estables y ampliamente adoptadas en la industria del software. La elección del stack buscó asegurar un equilibrio entre rendimiento, mantenibilidad, seguridad y disponibilidad de recursos a largo plazo. A continuación se presenta una descripción ampliada de cada uno de los componentes utilizados, explicando por qué fueron seleccionados y cómo contribuyen al funcionamiento global del sistema.

\subsubsection{Frontend: React + TypeScript}

React es actualmente una de las bibliotecas más populares para el desarrollo de interfaces web. Esto se debe a su forma de trabajar basada en componentes reutilizables, lo que permite construir pantallas modulares, organizadas y fáciles de mantener. Este enfoque facilita escalar la aplicación sin que la complejidad de la interfaz crezca de manera descontrolada.

Otro aspecto clave es su modelo de actualización mediante virtual DOM, que mejora considerablemente el rendimiento al renderizar solo lo necesario en cada cambio de estado. A esto se suma un ecosistema muy amplio de librerías, herramientas y documentación que acelera el desarrollo y permite incorporar funcionalidades avanzadas sin necesidad de reinventar soluciones existentes.

El uso de TypeScript complementa a React proporcionando tipado estático, ayudando a prevenir errores frecuentes en aplicaciones JavaScript puras y mejorando la claridad y mantenibilidad del código. Esta combinación hace que el frontend sea sólido, predecible y fácil de extender en el futuro.

\subsubsection{Backend: Java + Spring Boot}

Para el backend se optó por Java junto con Spring Boot, una de las combinaciones más confiables en entornos empresariales. Spring Boot permite crear aplicaciones robustas con una configuración mínima y una estructura bien organizada desde el inicio.

Entre sus principales ventajas para este proyecto se encuentran:

\begin{itemize}
    \item Integración nativa con Spring Security, que facilita la implementación de autenticación, autorización y manejo de permisos por roles.
    \item Soporte simple para APIs REST, ideal para la comunicación con el frontend.
    \item Inyección de dependencias, que fomenta una arquitectura modular y fácilmente testeable.
    \item Extensa documentación y comunidad activa, lo que asegura mantenimiento y soporte a largo plazo.
\end{itemize}

El uso de Java 21, una versión LTS (Long-Term Support), aporta estabilidad y asegura compatibilidad por varios años, permitiendo que el sistema se mantenga vigente sin necesidad de migraciones apresuradas.

\subsubsection{Base de Datos: PostgreSQL}

PostgreSQL fue seleccionada como base de datos por ser una de las opciones open-source más maduras y confiables disponibles en la actualidad. Es ampliamente utilizada en sistemas que requieren precisión, consistencia y soporte para operaciones complejas, características esenciales para manejar análisis, lotes y configuraciones del INIA.

Entre sus puntos fuertes destacan:

\begin{itemize}
    \item Soporte a transacciones ACID.
    \item Operaciones eficientes incluso con grandes volúmenes de datos.
    \item Flexibilidad para definir estructuras avanzadas, tipos personalizados y validaciones.
    \item Estabilidad comprobada en entornos de producción.
\end{itemize}

Estas características la convierten en una alternativa ideal para instituciones que dependen fuertemente de la integridad del dato.

\subsubsection{Infraestructura y Arquitectura}

La arquitectura elegida para el sistema sigue un modelo cliente-servidor, donde el frontend (React + TypeScript) interactúa con el backend (Spring Boot) mediante servicios REST. Esta separación clara de responsabilidades permite desarrollar y desplegar ambos componentes de manera independiente.

Además, el proyecto se divide en tres capas principales:

\begin{itemize}
    \item \textbf{Presentación:} la interfaz web que utilizan los usuarios.
    \item \textbf{Servicios:} lógica de negocio y validaciones.
    \item \textbf{Datos:} persistencia en la base PostgreSQL.
\end{itemize}

Este enfoque favorece la organización, la extensibilidad y el mantenimiento futuro del sistema.

\subsection{Soluciones Similares}

\subsubsection{Airtable}

Airtable combina conceptos de base de datos con la interfaz amigable de una hoja de cálculo. Es una plataforma low-code orientada a pequeños proyectos y equipos que necesitan digitalizar procesos sin desarrollar software propio.

\textbf{Ventajas:}
\begin{itemize}
    \item Interfaz simple y accesible
    \item Colaboración en tiempo real
    \item APIs integradas
    \item Automatizaciones básicas
\end{itemize}

\textbf{Limitaciones:}
\begin{itemize}
    \item No escala para flujos complejos
    \item Restricciones para reglas de negocio avanzadas
    \item Dependencia de licencias externas
\end{itemize}

\begin{figure}[H]
    \centering
    \includegraphics[width=0.9\textwidth]{images/airtable.png}
    \caption{Interfaz de Airtable}
    \label{fig:airtable}
\end{figure}

\subsubsection{Smartsheet}

Smartsheet ofrece una experiencia similar a Excel, pero con mayor control, trazabilidad y herramientas para flujos de trabajo.

\textbf{Ventajas:}
\begin{itemize}
    \item Gestión de proyectos y procesos
    \item Reportes avanzados
    \item Automatizaciones integradas
\end{itemize}

\textbf{Limitaciones:}
\begin{itemize}
    \item Alto costo según uso
    \item Menor flexibilidad frente a un desarrollo a medida
    \item Dependencia del ecosistema propietario
\end{itemize}

\begin{figure}[H]
    \centering
    \includegraphics[width=0.9\textwidth]{images/smartsheet.png}
    \caption{Interfaz de Smartsheet}
    \label{fig:smartsheet}
\end{figure}